\section{Results}
\label{sec:results}

All results come from the design of Sec.~\ref{sec:methods}, scored with the primary evaluator and, in parallel,
with the preregistered one. Unless stated otherwise, \textsf{C1}, \textsf{C1}$'$, \textsf{C0} and the $\alpha=0$
control come from one reference ensemble ($n=200$, $T=600$ weeks), \textsf{C2} from its production ensemble
($n=200$, $T=700$), and the long-lived control from the initial condition of \textsf{C1}.

\subsection{The envelope is generic; the rules add the classes}
\label{sec:res-null}

\textsf{C1} passes each applicable essential pattern in at least 99\% of its 200 runs and the primary outcome
O17p in $f=0.995$ [0.97, 1.00] (0.98 [0.95, 0.99] in an independent ensemble at $T=700$; 0.98 with the
preregistered evaluator). Table~\ref{tab:res-null} shows how much of this is evidence for the rules. Five of the
eight essential criteria, O4, O5, O6, O7 and O10, pass in at least two of the four null models and are generic:
they verify that the model has a boom, an end of natality and an extinction, which crowding damage and the
long-lived demography already give. Extinction in particular is not a result of the rules: the long-lived
control, the core model with the four rules switched off, $\alpha=1.4$ and $a_{\max}=120$ weeks (Sec.~\ref{sec:model}),
dies out in 0.98 of its runs, and in 199 of 200 as a null model run from the initial condition of \textsf{C1}
[Fig.~\ref{fig:res-rules}(a)]. The low peak with free space (O4) is generic too: without any rule, at
the calibrated attraction $\alpha=4$, the peak stays at $\rho_{\rm pk}=N_{\rm pk}/K_\gamma=0.45$ with $f_0=0.74$
of the cells empty. What the null models lack, and the rules add, is the coexistence of crowded and isolated
deteriorated classes (O13) and the persistence of the isolated class (O13p), which discriminate against three of
the four null models (both are present in $-$R2$^*$, where the classes form by full inheritance) and separate
\textsf{C1} from the partial ablations: \textsf{C0} passes every instantaneous pattern (O17 in 0.92) but its
isolated class is never persistent (O17p $=0$), and without attraction the colony fills 71\% of $K_\gamma$ with
3\% of empty cells and fails O4. O17p fails in every null model and in both partial ablations. The failure of the
crowded-born cohorts (O11) is a special case. Scored at six weeks, as in the primary evaluator, it discriminates
against every null model, but partly by definition: pups are born with $s\le w\,s_f\le0.2$, the threshold of the
criterion, and are still inside the plasticity window (with pups born at $s_{\rm base}=0.25$ instead of 0, O11
falls to 0.08). Re-scored at the end of the window (last row of Table~\ref{tab:res-null}), \textsf{C1} barely
changes (O11 0.995, O17p 0.990), but O11 then also passes in 0.81 [0.75, 0.86] of the long-lived runs and in 0.755
[0.69, 0.81] of the runs without rules, whose late cohorts even end lower than those of \textsf{C1} (median
$\bar s$ 0.074 and 0.055 against 0.130): their failure is accumulated crowding damage. The evidence of O17p against
those two null models therefore rests on O13 and O13p; against $-$R2$^*$ it rests on O11 at six weeks and is weaker
at $a_w$ (0.21).

\begin{table*}[t]
  \centering
  \caption{Which essential patterns carry evidence for the rules: fraction of runs that pass ($n=200$ per model,
  primary evaluator, NA counted as a failure) in \textsf{C1} and in the four null models run from its initial
  condition: no crowding damage ($\gamma=0$); full inheritance without social learning ($-$R2$^*$: $w=1$,
  individual recovery at $r=0.01$); long-lived demography without new rules (long); \textsf{C1} without its four
  rules (none). ``Produced by'': the ingredient that makes the pattern pass. Generic: passes ($f\ge0.8$) in at
  least two null models. Discriminates: upper Wilson bound below 0.8 and Newcombe interval of the difference with
  \textsf{C1} excluding zero. $^\ddagger$Holds by design. O6 and O7 share a row (same pass rates). The last row
  re-scores O11 on the same runs with the late cohorts measured at the end of the plasticity window ($a_w=12$
  weeks) instead of at six weeks. Desirable patterns, partial ablations and preregistered criteria: SM, Tables~S4--S6.}
  \label{tab:res-null}
  \footnotesize
  \setlength{\tabcolsep}{2.7pt}
  \begin{tabular}{@{}llcccccl@{}}
    \toprule
    Pattern & Produced by & \textsf{C1} & $\gamma{=}0$ & $-$R2$^*$ & long & none & Evidence \\
    \midrule
    O4 low peak with free space & $\alpha=4$ with crowding damage & 1.00 & 0.00 & 0.99 & 0.17 & 1.00 & generic \\
    O5 natality ends at high $N$ & demography & 0.99 & 1.00 & 0.97 & 0.89 & 0.91 & generic \\
    O6, O7 no rebound; extinction & crowding damage, senescence & 1.00 & 0.00 & 1.00 & 0.99 & 1.00 & generic \\
    O10 adult deterioration before the peak & crowding damage & 1.00 & 0.00 & 1.00 & 1.00 & 1.00 & generic \\
    O11 late cohorts fail (age 6 weeks) & R2 with $s_{\rm base}=0$ & 1.00 & 0.00 & 0.00 & 0.00 & 0.00 & discriminates, 4 of 4 \\
    O13 two deteriorated classes & AV or R7 & 1.00 & 0.00 & 1.00 & 0.00 & 0.00 & discriminates, 3 of 4 \\
    O13p persistent isolated class$^\ddagger$ & R7 ($c_R=2$) & 1.00 & 0.00 & 1.00 & 0.00 & 0.00 & discriminates, 3 of 4 \\
    \midrule
    O17p primary outcome & all of the above & 0.995 & 0.00 & 0.00 & 0.00 & 0.00 & discriminates, 4 of 4 \\
    \midrule
    O11 scored at age $a_w$ & crowding damage, mostly & 0.995 & 0.00 & 0.21 & 0.81 & 0.755 & not generic; discriminates, 2 of 4 \\
    \bottomrule
  \end{tabular}
\end{table*}

\paragraph{Ablations and the minimal candidate.}
Each single ablation of \textsf{C1} fails in a specific way [Fig.~\ref{fig:res-rules}(b)]: without R7 the
isolated class is never persistent (O17p $=0$, O17 in 0.94); without R2 the late cohorts do not fail at six weeks
(O17p $=0.00$ [0, 0.03]); without R1 adults recover when the colony thins, natality resumes, only 23\% of the
colonies are extinct by week 600 and O17p is 0.21 [0.15, 0.29]; without AV O17p falls to 0.79 [0.71, 0.85],
through O5 alone. Every multiple ablation has O17p $=0$. Inheritance is not needed ($w=0$ with $r=0.15$ gives
1.00), so social learning is the essential part of R2; R7 needs its preference ($\pi=0$ gives 0) and, under
synchronous moves, its hard capacity (a soft capacity gives 0, because simultaneous moves send many deteriorated
animals into the same refuge; under a random sequential update it gives 0.995 against 1.00). The preregistered
intersection--union test of Q2 (``each rule is necessary'') is decided by AV: with the primary evaluator
$p_{\rm Holm}=2.2\times10^{-7}$, rejected; with the preregistered evaluator $f(-{\rm AV})=0.90$ and
$p_{\rm Holm}=0.067$, not rejected. The rejection is thus a consequence of the primary O5 alone, a generic
criterion, and with either evaluator the preregistered prediction $f(\textsf{C1}-{\rm AV})\le0.5$ fails (0.79): the
necessity of AV, as defined in the plan, is refuted. AV contributes by shortening the tail of births ($\tau_B=6.9$
against 8.7 weeks), supplying the senescent decline (O8 in 0.62 against 0.21) and \emph{reducing} the empty space
at the peak ($f_0=0.16$ against 0.52), which is present without any rule. The minimal set for the primary outcome
is therefore R1 + R2 + R7, the candidate \textsf{C1}$'$ (O17p 0.835 [0.78, 0.88], failing O5 in 15\% of the runs),
in which each of R1, R2 and R7 remains necessary (O17p of 0 without each, $n=40$). This minimality is conditional
on the evaluation thresholds: O17p of \textsf{C1}$'$ is 0.83 for class thresholds $\theta\le0.125$ and 0.33 at
$\theta=0.15$, while \textsf{C1} holds 0.96 at 0.15 (SM, Sec.~S3).

\begin{figure}[t]
  \centering
  \includegraphics[width=\columnwidth]{fig_pre_rules.pdf}
  \caption{What the rules change. \textbf{(a)}~Kaplan--Meier survival of the colonies for the candidates
  ($n=200$; $T=600$, or 700 for \textsf{C2}) and the controls (the core model without the four rules and $\alpha=1.4$: short-lived,
  $a_{\max}=60$; long-lived, $a_{\max}=120$; $n=100$, $T=1000$), with log--log Greenwood bands for \textsf{C1} and
  the short-lived control. Every model with the long-lived demography dies out within a narrow window; only the
  short-lived control survives and oscillates (53 of 100 colonies alive at $T$). \textbf{(b)}~O17p of \textsf{C1}, its rule ablations
  and implementation variants with the primary evaluator (filled, 95\% Wilson interval) and the preregistered
  evaluator (open); single ablations and variants $n=120$, multiple ablations and $\gamma=0$ $n=40$. The right
  margin names the essential criterion with the lowest pass rate where O17p $<0.95$; the dotted line is the 0.8
  threshold.}
  \label{fig:res-rules}
\end{figure}

\subsection{Trajectories}
\label{sec:res-trajectories}

The time course of \textsf{C1} (Fig.~\ref{fig:res-trajectories}) is a single boom, a plateau and an extinction
by old age. The mean social state of the adults (age $\ge6$ weeks) falls below 0.5 at week 13 [13, 14], while the
population is still growing; part of this fall is dilution by pups born near $s=0$, and measured on the animals
past the plasticity window or on the founders the crossing comes at week 16, so deterioration takes of the order of
one generation, $\Gamma=t_{\rm det}/T_{\rm gen}\approx1.1$ ($t_{\rm det}\approx10$ weeks, $T_{\rm gen}\approx9$).
The socially functional growth phase is short either way (O10b in 0.02 of the runs with the adult mean, 0.28 on the
animals past the window): a phase B of 13 weeks that holds two thirds of the births, against 30 weeks in
Universe~25. The population comes within 2\% of its maximum at $t'_{\rm pk}=34$ [32, 36] weeks, 99\% of the births
have occurred by week 40, the last surviving birth is at week 59, the plateau lasts 51 [49, 53] weeks, and the
colony then ages at one week per week and dies out at week 169 [167, 171] (Kaplan--Meier median). The peak uses
$\rho_{\rm pk}=0.36$ of the damage-threshold capacity and leaves 16\% of the cells empty (Universe~25: 0.57 and
about 20\%; an agreement in sign only, since $f_0$ was used in the calibration), and 17\% of the adults form a
persistent isolated class (0.6\% in \textsf{C0}). The desirable patterns are only partly reproduced: O8 passes in
0.62 [0.56, 0.69] of the runs, O9 in 0.45 and O13b in none (SM, Table~S3). Without mortality before senescence the
maximum $t_{\rm pk}=\arg\max N$, at week 55 [50, 60], falls on the last birth before it in 69\% of the runs, 21
weeks after the plateau begins; the decline lasts $(T_{\rm ext}-t_{\rm pk})/t_{\rm pk}=2.07$ peak times with
$\arg\max N$, close to Calhoun's 1.84--2.2, but 3.97 with $t'_{\rm pk}$, so this agreement is not a robust result.
The short-lived control, the core model with the four rules switched off, $\alpha=1.4$ and a short life
($a_{\max}=60$, $a_{\rm rep}=50$ weeks), fails O11 and O13 in every run and 53\% of its colonies oscillate. \textsf{C2}, which
ties the retreat to maternal socialization, makes its isolated class a later-born cohort (O13b at least marginal in
200/200 runs against 0/200 in \textsf{C1}) but pays with its demography: O8 fails in every run, O5 falls to 0.815
and O17p to 0.705 [0.64, 0.76], against 0.98 in the \textsf{C1} ensemble of the same stage (Holm
$p=4\times10^{-14}$ over the 29 rows of the \textsf{C1}--\textsf{C2} family), whereas the calibration had predicted
no difference (SM, Table~S7).

\begin{figure*}[t]
  \centering
  \includegraphics[width=\textwidth]{fig_pre_trajectories.pdf}
  \caption{Time course of the reference ensemble ($n=200$ each). \textbf{(a)}~Population relative to the
  damage-threshold capacity $K_\gamma=m_c|G|$: median and 10--90\% band for \textsf{C1} (thin lines are 12
  individual runs), and medians for \textsf{C1}$'$, \textsf{C0} and the $\alpha=0$ control. \textbf{(b)}~Weekly
  births (median). \textbf{(c)}~Mean social state of the adults (age $\ge6$ weeks; lines as in a), median over
  living colonies while at least half of them are alive; for \textsf{C1} also the mean over the animals past the
  plasticity window (dotted) and over the founders (dashed), which are not diluted by the newborn pups. Grey dotted
  lines: the 0.5 level of O10 and the 0.2 deterioration threshold. \textbf{(d)}~Mean age, with a reference line of
  slope one: after natality stops, the colony simply ages.}
  \label{fig:res-trajectories}
\end{figure*}

\subsection{Irreversibility: R1 and the transplants}
\label{sec:res-irreversibility}

Calhoun reports that small groups moved from mid-phase D to uncrowded universes failed to reproduce, even with
partners raised without crowding \cite{Calhoun1973,Marsden1972}. The model reproduces this (O12), but by
construction, so the transplants are a consistency check of R1 rather than a finding (Fig.~\ref{fig:res-transplant}).
Four males and four females of each \textsf{C1} colony, taken at $1.75\,t'_{\rm pk}$ (week 60; ages 37--64 weeks),
never found a colony in a \textsf{C1} enclosure (0/60), nor with healthy partners (0/120), but do in 51/60 trials
when the same animals are placed in an enclosure without R1 and R2 with individual recovery at $r=0.1$ (arm $X$;
exact McNemar $p=4\times10^{-16}$), so the preregistered prediction for arm $X$ ($\ge0.25$) holds (Q1). These
animals have $s\approx0$ (median 0.000, maximum 0.36), and fecundity scales as $s_f^2s_m$, so under R1 they cannot
found a colony with any partner. The rate in arm $X$ is a matter of age: at the preregistered sampling time
$1.75\,t_{\rm pk}$ (week 97; ages 59--80, with the fertile window closing at 80) the same colonies give 12/60, and
at week 95 in every arm, the common time that removes the confound between rule and sampling time, 5/60. A new arm
samples adults past the plasticity window with $0.2\le s\le0.5$ (median 0.36; weeks 8--11): under R1 they found a
colony in 5/60 trials, against 57/60 for the same animals in arm $X$ and 50/60 without R1, and the adult
descendants of the colonies they found settle at the social state of their parents (0.28--0.45 in \textsf{C1}).
The irreversibility of R1 + R2 is thus strict for $s\approx0$; for intermediate animals R1 acts through fecundity,
not through a collapse of the lineage, which confirms only the first part of the mean-field lineage argument of
\extended\ (SM, Sec.~S5).

\begin{figure}[t]
  \centering
  \includegraphics[width=\columnwidth]{fig_pre_transplant.pdf}
  \caption{Transplant experiments (60 replicates per arm and phase). $P_T$ is the probability that four males and
  four females moved to an empty lattice reach 80 individuals within 150 weeks. \textbf{(a)}~Animals taken in
  phase B and in phase D at three sampling times: $1.75\,t'_{\rm pk}$ of each colony, $1.75\,t_{\rm pk}$ with
  $t_{\rm pk}=\arg\max N$ (preregistered protocol), and week 95 in every arm; dotted lines mark the O12 thresholds.
  \textbf{(b)}~Adults past the plasticity window with $0.2\le s\le0.5$, alone and paired with healthy animals of
  the other sex. Arm $X$ moves the animals of the \textsf{C1} arm, the same individuals, into an enclosure without
  R1 and R2 with individual recovery at $r=0.1$.}
  \label{fig:res-transplant}
\end{figure}

\subsection{A boom of one to two generations}
\label{sec:res-boom}

Following every animal of the reference ensemble by generation shows how shallow the collapse is. The founders
deteriorate by crowding during their own adult life (they leave the plasticity window with $s=0.90$ but have
$s=0.20$ at 26 weeks of age) and produce 69\% of all births, about 6\% of their potential fecundity. Their
offspring are born at $s_{\rm nb}=w\,s_f\simeq0.19$, below the deterioration threshold, and gain only $+0.04$ by
the end of the window, because their teachers are already deteriorated. The net reproductive number falls by a
factor of about ten in a single step, from 3.8 daughters per founder female to 0.40 in the first generation born,
then to 0.11 and 0.05; the second generation adds 27.6\% of the births and the third and later ones 3.2\%. The
mean generation of the births is 1.34 in \textsf{C1}, against 1.95 [1.90, 1.99] in the long-lived control and 1.59
in \textsf{C1} without its rules, whose first generation still has $R_{\rm net}=1.21$ and 0.66: the rules, R2 in
particular, shorten the generational depth, because pups no longer inherit the competence of the few healthy
mothers. R2 does not produce a cascade over many generations; it makes the failure of the first generations born
during the boom irreversible for the colony, in every variant of Sec.~\ref{sec:res-robustness}. Births are
front-loaded: half of them occur before $0.36\,t'_{\rm pk}$ and 93\% before $0.8\,t'_{\rm pk}$, the start of the
window of the strict O13b criterion, which therefore cannot pass. In \textsf{C1} the isolated and the crowded
classes are born together (median birth weeks 10 and 11), so the persistent isolated class is made of refugees of
the early boom, most of whom entered the refuges as pups, not of Calhoun's late-born ``beautiful ones'';
\textsf{C2} obtains the later-born class (median birth week 22 against 11) at the demographic cost described
above. The simulated dynamics is a boom of one or two generations rather than Calhoun's multi-generational
sequence (Sec.~\ref{sec:discussion}).

\subsection{Robustness and the central limitation}
\label{sec:res-robustness}

Fig.~\ref{fig:res-rob} reports \textsf{C1} and \textsf{C1}$'$ under the robustness variants ($n=200$ each); by the rule fixed before the runs, a variant preserves the result if O17p $\ge0.8$ (SM,
Table~S9). With a constant background hazard $\mu_{\rm bg}$ at all ages the plateau shrinks (51 weeks at 0, 35 at
0.001, 17 at 0.003 and 10 at 0.008 per week) and O17p falls to 0.83, 0.63, 0.55 [0.48, 0.62] and 0.06 in
\textsf{C1} (0.67, 0.38, 0.28 and 0.04 in \textsf{C1}$'$). The only criterion that fails is O5, and it fails
through its timing clause, not through its level: at $\mu_{\rm bg}=0.003$ O5 is MARGINAL in 83 of the 200 runs
and FAIL in 7, because the start of the plateau moves from week 34 to 29 while $t_{B99}$ moves from 40 to 43, so
$(t_{B99}-t'_{\rm pk})/t'_{\rm pk}$ crosses the 0.5 threshold of O5 while natality still ends with the population
near its maximum. The classes, O11 and the low peak are unchanged, and O17p restricted to the criteria that
discriminate against the null models (O11 at six weeks, O13, O13p and the anti-patterns) is 1.00 at every
$\mu_{\rm bg}$. With initial ages drawn uniformly from 4--52 weeks instead of all equal, O17p is 0.93 [0.89, 0.96];
from 4--80 weeks, which includes post-reproductive founders, 0.63 [0.56, 0.69] (0.68 and 0.31 in \textsf{C1}$'$),
again through the timing clause of O5 (73 MARGINAL and 1 FAIL of 200) and with the restricted outcome at 1.00. The
timing of the end of natality relative to the start of the plateau is therefore a property of the synchronous
initial condition and of the absence of mortality before senescence, not of the rules; the end of natality near
the maximum itself persists. Since O5 is generic, these variants remove the generic part of the primary outcome,
not the part that carries evidence for the rules; but the primary outcome \emph{as defined} holds only without,
or with very weak, background mortality, which is the central limitation of the result
(Sec.~\ref{sec:discussion}). The other variants preserve it: a nest period of three weeks gives 0.835 [0.78, 0.88],
borderline because the interval contains 0.8 (0.60 in \textsf{C1}$'$); a random sequential update gives 1.00 with
the hard and 0.995 with the soft refuge capacity; and the exact refuge capacity gives 0.995 against 0.995 with a
single-pass admission. Pups born above the deterioration threshold ($s_{\rm base}=0.25$ or 0.5) give O17p
of 0.08 and 0.00, entirely through O11 scored at six weeks; scored at the age $a_w$ the same runs give O17p of
0.965 and 0.76, because by the end of the window the late cohorts are below 0.2 whatever their birth state. Over
lattice sides $L=15$--60 at fixed density $N_0/L^2=0.44$ ($n=60$; SM, Table~S10), the peak density (0.35--0.37)
and the persistent fraction (0.17--0.18) do not change and O17p rises from 0.82 [0.70, 0.89] at $L=15$ to 1.00 at
$L\ge45$; the spread of the extinction time falls from 0.15 to 0.08 together with the relative spread of the last
birth (0.58 to 0.18), while $T_{\rm ext}-t_{\rm last}$ stays at 110--112 weeks [Fig.~\ref{fig:res-rob}(c)]: the
spread of $T_{\rm ext}$ is that of the end of natality, and a narrow spread is not a signature of the rules (0.045
without attraction, 0.160 in the long-lived control). O17p reaches 0.8 only for seeding densities
$0.2\lesssim N_0/L^2\lesssim0.9$, and with Calhoun's four founding pairs no run is Calhoun-like (the best point of
an exploratory founder plane reproduces O1--O3 together with O17p in 0.43 [0.27, 0.61] of the runs). Around
\textsf{C1} the primary outcome holds over part of the explored planes (30 runs
per cell): O17p is at least 0.83 for $w\le0.3$ and $r\le0.10$ and falls at $w\ge0.5$ and $r\ge0.2$ through O11;
every finite window $a_w\le24$ gives certain extinction; in the $(\alpha,\kappa)$ plane O17p is at least 0.8 at 34
of the 48 points with $\alpha\ge2$, failing where the free space of O4 is lost; and a persistent isolated class
requires a strong refuge preference and a total refuge capacity $f_Rc_R\ge0.4$.

\begin{figure*}[t]
  \centering
  \includegraphics[width=\textwidth]{fig_pre_robustness.pdf}
  \caption{Robustness of the primary outcome ($n=200$ per arm unless stated). \textbf{(a)}~O17p against the
  background mortality $\mu_{\rm bg}$ for \textsf{C1} and \textsf{C1}$'$ ($n=60$ at the intermediate values), with
  the primary (filled) and the preregistered evaluator (open); dotted line: the 0.8 threshold. \textbf{(b)}~Variants
  of \textsf{C1} and \textsf{C1}$'$ (vertical lines: the reference ensembles and the 0.8 threshold).
  \textbf{(c)}~Lattice side at fixed seeding density ($n=60$ per size): Kaplan--Meier spread of the extinction
  time (interquartile range over median), relative spread of the last birth, and O17p.}
  \label{fig:res-rob}
\end{figure*}

\subsection{Feasible region and sensitivity}
\label{sec:res-global}

Over two independent 15-factor Latin hypercubes (400 points, 10 runs each, plus 40 runs at 18 boundary points),
2.0\% [0.7, 3.6] of the points are Calhoun-like with the primary evaluator ($k/n\ge0.8$; bootstrap over points),
the mean of O17p over the box is 4.0\% [2.6, 5.9], and an empirical-Bayes fraction that shrinks each point towards
a logistic emulator is 1.8\% [0.7, 3.2]; with the preregistered evaluator the same runs give 7.2\%, 11.0\% and 6.9\%
[4.8, 9.2]. Every such fraction is relative to the box and to the sampling scales of Table~\ref{tab:params}. Part of
the small fraction is definitional and part is mechanism: with $c_R=3$ refuge residents exceed the isolation
threshold and none of the 100 such points is Calhoun-like; restricting successively to $c_R\le2$, $m_c\le10$,
$f_Rc_R\ge0.3$ and $\pi\ge100$ (chosen after the runs) raises the fraction to 2.7\%, 3.1\%, 5.6\% and 10.8\%
[2.7, 21.6], and even then 33 of the 37 points fail, through O11 (16), the persistence of the class (14), O5 (9)
and O4 (9). The Sobol' decomposition (11 factors; SM, Fig.~S1) finds O17p equal to zero at 78\% of its points; its
total indices are led by $c_R$ ($0.49\pm0.16$), $w$ (0.43), $\kappa$ (0.41), $r$ (0.32) and $m_c$ (0.31), which
cannot be ordered, and its first-order indices add up to only 0.15, so it is almost entirely interactions, while
the continuous outputs are nearly additive; the window $a_w$ does not exceed its noise floor.

\subsection{Confirmatory tests}
\label{sec:res-Q}

With the primary evaluator six of the seven preregistered null hypotheses are rejected after Holm adjustment and
the quantitative predictions are met for Q1, Q3, Q4 and Q7 and not for Q2, Q5 and Q6; with the preregistered
evaluator five are rejected and four predictions hold (Q3, Q4, Q6 and Q7). The two counts differ in three places:
Q2 is rejected only with the primary O5, the prediction of Q1 holds only at the earlier transplant time, and the
prediction of Q6 holds only with the preregistered O11 (SM, Table~S7). Of the predictions that hold, Q1 is a
consistency check of R1 and Q4 of the demography: once births stop, extinction is the death of the last survivor
of a closed, ageing population, so the expected $T_{\rm ext}-t_{\rm last}$ grows one-to-one with $a_{\max}$ and
only logarithmically with the size of the last cohort (slope $0.996\pm0.014$ over 300 runs). Q3 confirms a
near-deterministic collapse against the short-lived control (200/200 extinct by $T=600$ against 47/100 by $T=1000$; coefficient
of variation of $T_{\rm ext}$ 0.066 [0.059, 0.071]), but the long-lived demography alone goes extinct in 0.98--1.00
of the runs ($p=0.11$), so extinction does not discriminate. Q5, a coexistence band of O13 centred at the
preferred occupancy $n^*=1/\kappa-1/\alpha\approx m_c$, is refuted in the opposite direction (O13 passes in at
least 95\% of the runs at 46 of the 49 points with $n^*>0.5$): with refuges the isolated class no longer depends on
the balance between attraction and aversion. Q6, that the window weighs less than the learning rate ($0<\theta<1$
in $\log r+\theta\log a_w$), is not supported with the primary evaluator: the Calhoun-like boundary of the
$(r,a_w)$ plane is nearly vertical in $r$ and the single-index fit gives $\hat\theta=-0.15$ [$-0.20$, $-0.05$] in a
model that an additive one beats ($\Delta{\rm QAIC}=28$); with the preregistered evaluator it gives $\hat\theta=0.30$
[0.15, 0.50]. Q7 confirms an endogenous collapse without crowding damage: at $\gamma=0$ the four subcritical points
($w\le0.1$, $r\le0.03$) went extinct in 120/120 runs while all 30 \textsf{C1} colonies exploded, but the colonies
of the corner first overshoot the capacity ($\rho_{\rm pk}=0.91$--1.60), so none is Calhoun-like; with the window
varied, the transition in $\Lambda=r\,a_w$ falls in the same grid bracket for $a_w=12$ and 24 and one step lower
for $a_w=6$, and the preregistered logistic model is separated (51 of 54 points have $P_{\rm ext}\in\{0,1\}$), so no
single exponent is reported as inference (SM, Sec.~S4). The mean-field analysis of \extended\ mostly describes the
short-lived control (a superstable limit cycle sustained by healthy pockets of uncrowded individuals, without
secular drift of the peaks); what carries over to \textsf{C1} is the phase-D law of Q4, the capacity bound of the
persistent class $\Phi^{\rm pers}_{\rm BO}\le f_Rc_R|G|/N_{\rm pk}$, and the heuristics behind Q6 and Q7, which the
data support only in part.
